\section{From workload to enterprise economics}
\label{sec:economics}

\subsection{Comparable resource costs}

Let $w_s$ be a human-hour valuation, $a_s$ variable nonlabor cost per visit, $F_s$ fixed nonlabor
cost, and $\ell_s$ expected attributed error loss per visit under regime $s$. A common-boundary
average cost is
\begin{equation}
\bar c_s = w_s\bar h_s + a_s + \frac{w_sB_s+F_s}{\lambda_s} + \ell_s. \label{eq:avgcost}
\end{equation}
Human upkeep is included through $B_s$, not charged again in $F_s$. Both regimes bear error losses.
Different $w_s$ values allow a more highly paid specialist residual. For mixed roles, the labor
valuation becomes a sum of hours times role-specific costs.

For the IT Production and Standards population under S2 ($h_0=30$ hours, $q=0.24$), $\bar h_A=17.18$
and $B_A/\lambda=1.2$ hours. Take artificial monetary
units with $w_H=80$, $w_A=100$, $a_H=20$, $a_A=50$, $F_H=1000$, $F_A=4000$, and $\ell_H=30$,
$\ell_A=50$ at 100 monthly visits. Then
\begin{equation}
\bar c_H = 2460.00, \qquad \bar c_A = 1978.00. \label{eq:costexample}
\end{equation}
Holding the other terms fixed, $\ell_A=532.00$ is the error-loss value at which the costs equalize. A
separate security or safety constraint can still prohibit deployment below this financial threshold.
The numbers illustrate opportunity costs, not market salaries, purchase prices, or payroll reductions.

\subsection{Average cost, scale, and system return}

When variable cost $v_s=w_s\bar h_s+a_s+\ell_s$ is constant and fixed cost is $J_s=w_sB_s+F_s$, total
period cost is $T_s(\lambda)=J_s+\lambda v_s$. Average cost is $J_s/\lambda+v_s$ and marginal cost is
$v_s$. Fixed assurance becomes cheaper per case as volume grows, but capacity thresholds or
congestion can change this relationship. A high-volume standardized route may justify an investment
that is uneconomic for a rare project type; Section~\ref{sec:routes} gives the workload counterpart.

At pipeline level, a shared incident must be counted once. An aggregate comparison is
\begin{equation}
C_s=\sum_i\bigl[w_{si}L_{si}+\lambda_{si}a_{si}+F_{si}\bigr]+\mathbb{E}[\mathcal{L}_s], \label{eq:pipeline}
\end{equation}
where $\mathcal{L}_s$ is total process loss over the period. Shared labor and technology are
allocated once; local losses are not added again to the aggregate loss. The comparison covers the
same requested business outputs and includes any different numbers of visits required to produce them.

For an investment horizon $0,\dots,T$, let $V_t$ be nonduplicated attributable monetary benefits,
$I_t$ incremental implementation and operating expenditure, and $d$ the discount rate. One explicit
ROI convention is
\begin{equation}
\mathrm{ROI}=\frac{\sum_{t=0}^{T}(V_t-I_t)/(1+d)^t}{\sum_{t=0}^{T}I_t/(1+d)^t},
\qquad \sum_{t=0}^{T}I_t/(1+d)^t>0. \label{eq:roi}
\end{equation}
Benefits already net of a cost cannot subtract that same cost again. Resource savings and cash
savings require distinct benefit series when salaries remain committed. Transition work, migration,
training, and any redundancy costs belong to the investment comparison.

Upstream gains can reach several departments, but their returns are not automatically superadditive.
If $B(S)=C(\varnothing)-C(S)$ is net saving for intervention set $S$, the interaction
\begin{equation}
B(\{u,v\})-B(\{u\})-B(\{v\}) \label{eq:interaction}
\end{equation}
is positive only when the joint intervention saves more than the isolated gains combined. Shared
information, integration costs, or overlapping work can make the interaction negative. This test is
the economic basis for investigating upstream leverage, rather than assuming it from process order.
